% Section 5, Testing the model (split from the quantification section, Oct 25 2026).
% 2007 fit: output/model/calibration/2026-10-19_identification_newacct/ (moments/center.json) and
% output/model/calibration/2026-10-21_lifecycle_newacct/. 2023: experiments/housing_supply/20261021_free_entry/
% moments_2023_fi/ (free_entry_eta0.63 column; replay of the converged FIN063fi_fw history).

\section{Testing the model}
\label{sec:quant-fit}

In this section I compare the model with features of the data that the
calibration does not target. I focus on the margins that matter for the
policy experiments of Section~\ref{sec:policy}: how many children women have
and when, how households change their housing around a birth, and how
ownership and housing evolve over the life cycle.
Table~\ref{tab:quant-untargeted} reports the untargeted moments, and
Figures~\ref{fig:quant-family-size} to~\ref{fig:quant-event-scatter} compare
the full distributions and profiles. Appendix~\ref{sec:quant-2023} compares
the 2023 economy, the starting point of the policy experiments, with the
2023 data.

\begin{table}[t]
\centering
\caption{Moments not targeted, 2007}
\label{tab:quant-untargeted}
\small
\begin{tabularx}{\textwidth}{@{}X l r r@{}}
\toprule
Moment & Source & Data & Model \\
\midrule
First births at age 30+ (\%) & NCHS 2003--2006 & 24.9 & 24.2 \\
Rooms: 3+ versus 1--2 resident children & ACS 2005--2006 & 0.385 & 0.574 \\
First-birth ownership response & PSID event study & 0.261 & 0.141 \\
Wealth/income p90 / median, ages 76--84 & PSID & 3.52 & 3.27 \\
\bottomrule
\end{tabularx}
\par\smallskip
\begin{minipage}{0.95\linewidth}
\footnotesize\textit{Notes:} This table reports moments that do not enter
the loss \eqref{eq:q-loss}. The ownership response is the event-study
coefficient of Section~\ref{sec:empirical-evidence} on the model's
observation support, measured on simulated households with the same design.
% [MODEL] moments/center.json target_fit validation rows: share30 0.2421,
% family_rooms 0.5740, old_dispersion 3.2730 (7.9: 3.4656, same row); ownership
% event birth_dated Wp3 0.1410.
\end{minipage}
\end{table}

\subsection{Fertility}
The model reproduces the distribution of completed family size
(Figure~\ref{fig:quant-family-size}). The calibration targets the share of
women with no children and the share of mothers with one child. It does not
target how the remaining mothers split between two children and three or
more. In the model, 32.2\% of women aged 40 to 44 have two children and
30.4\% have three or more, against 34.4\% and 28.6\% in the CPS. The timing
of first births is also close to the data: 24.2\% of first births occur at
30 or later, against 24.9\% in the natality files. The right panel of
Figure~\ref{fig:quant-lifecycle-2007} shows the whole distribution. The
model places fewer first births between 26 and 33, 28.5\% against 33.6\%,
and more at 18 to 21.
% [MODEL] family size: kappa_identification/stage3/runs/center/.../observers.json uniform_birth_time
% parity_shares_40_44 (20.0/17.3/32.2/30.4); [DATA] 2026-10-25_data_pictures/cps_family_size_40_44.json
% (CPS June 2004/06, n = 10,872: 19.8/17.1/34.4/28.6).

\begin{figure}[tbp]
\centering
\includegraphics[width=0.55\linewidth]{../../output/model/calibration/2026-10-25_data_pictures/family_size_2007.pdf}
\caption{Completed family size, model and data}
\label{fig:quant-family-size}
\par\smallskip
\begin{minipage}{0.92\linewidth}
\footnotesize\textit{Notes:} Share of women aged 40 to 44 by number of
children ever born. Model: 2007 steady state. Data: June Fertility Supplement
of the CPS, 2004 and 2006 pooled.
\end{minipage}
\end{figure}

\begin{figure}[tbp]
\centering
\includegraphics[width=\linewidth]{../../output/model/calibration/2026-10-21_lifecycle_newacct/lifecycle_2007_newacct.pdf}
\caption{Life-cycle profiles in 2007, model and data}
\label{fig:quant-lifecycle-2007}
\par\smallskip
\begin{minipage}{0.92\linewidth}
\footnotesize\textit{Notes:} Homeownership and rooms by age of the household
head, and the distribution of first births by age of the mother. Model: 2007
steady state. Data: ownership and rooms (at most nine) from the 2005--2006
ACS; first births by four-year age group from the NCHS natality files for
2005--2007.
\end{minipage}
\end{figure}

\subsection{Housing around the first birth}
Figure~\ref{fig:quant-event-scatter} compares the change in rooms and in
ownership around the first birth in the model and in the PSID, window by
window. The calibration targets only the change in rooms three to four years
after the birth. The rest of the room response after the birth is close to
the data for about eight years: in the window of the birth, households add
0.78 rooms in the model and 0.62 in the data. Two differences stand out.
First, model households already add rooms in the years before the birth,
while in the data rooms are flat before it. Second, ownership rises around
the first birth by 14 percentage points in the model, against 26 in the
data.
% [MODEL] center.json first_birth_event_path, birth_dated (omitted -3/-2): <=-8 -0.65,
% -7/-6 -0.53, -5/-4 -0.20, -1/0 0.78, +1/+2 1.51, +3/+4 1.48, >=11 0.65; ownership_event Wp3 0.141.
% [DATA] rooms_coefficients.csv (A2h, common support), ownership A2h_own: Wp3 0.261.

\begin{figure}[tbp]
\centering
\includegraphics[width=0.85\linewidth]{../../output/model/calibration/2026-10-25_data_pictures/untargeted_event_scatter.pdf}
\caption{Housing around the first birth, model and data}
\label{fig:quant-event-scatter}
\par\smallskip
\begin{minipage}{0.92\linewidth}
\footnotesize\textit{Notes:} Each point is a two-year window around the
first birth, relative to the window two to three years before it. Data: the
PSID event study of Section~\ref{sec:empirical-evidence}. Model: the same
regression on simulated households in the 2007 steady state. The hollow
marker is the targeted room response.
\end{minipage}
\end{figure}

\subsection{Life-cycle profiles}
Figure~\ref{fig:quant-lifecycle-2007} compares ownership and rooms by age of
the head with the 2005--2006 ACS. The model reproduces the rise of
ownership with age, and between 34 and 49 the two profiles are within four
percentage points. It also reproduces the hump in housing size: rooms rise
until the forties and then decline, although the model profile is steeper.
Young households own too often in the model, 34\% at 18--21 and 44\% at
22--25, against 12\% and 24\% in the data. Older households keep their homes
in the model, since it has no health or long-term-care risk that would lead
them to sell.
% [MODEL] output/model/calibration/2026-10-21_lifecycle_newacct/lifecycle_2007_newacct.json;
% [DATA] output/model/production/2007/plotted_long.csv (ACS 2005/06, rooms capped at 9;
% NCHS 2005-07 first births).

\subsection{Wealth}
The model accounts for the dispersion of wealth among the old: the ratio of
the 90th percentile to the median of wealth relative to income at ages 76 to
84 is 3.27, against 3.52 in the PSID.
